\chapter{Conclusion and Future Work}

\section{Summary of the Thesis}
This thesis studied multimodal target trajectory completion, fusion, and behavior analysis under incomplete and heterogeneous observations. The core problem is that different sensing sources may observe the same physical target with different coordinate availability, sampling density, attribute form, and reliability. Directly combining such observations can easily produce unstable identities or noisy behavior results. To address this issue, the thesis built a progressive framework from optical reference retrieval to multimodal entity fusion and behavior analysis.

First, the thesis established a unified target-centered representation for heterogeneous observations. Optical, electromagnetic, SAR, thermal, and trajectory-based sources are organized with position, motion, attributes, confidence, and validity information. This representation provides the basis for subsequent module connection.

Second, for coordinate-missing optical images, the thesis adopted TF-VPR as the optical reference retrieval module. The module retrieves coordinate-known reference images instead of directly estimating geographic coordinates from a single query image. The retrieved Top-\(k\) references provide scene-level anchors for downstream target association, allowing coordinate-missing optical observations to enter the trajectory system.

Third, the thesis formulated multimodal target fusion as an entity-level registration problem. Completed trajectories from different modalities are encoded into entity embeddings, refined through multi-set graph matching, and merged through reliability-aware fusion. The result is a fused target representation that combines trajectory information, modality-specific attributes, and confidence information.

Finally, the fused entities were used for behavior analysis. Individual behavior is described through route, speed, and shape features, while group behavior is modeled through dynamic spatiotemporal relations. Individual and group evidence are then combined into anomaly scores, event decisions, and regional situation heatmaps.

Overall, the thesis provides a coherent technical route from incomplete multimodal observations to fused target representation and behavior-level interpretation. The main value of the work lies in connecting optical reference retrieval, entity-level fusion, and behavior analysis into one consistent framework.

\section{Limitations}
Several limitations remain. First, the optical completion stage still depends on the quality of both reference retrieval and downstream target association. TF-VPR provides scene-level candidate references, but target-level correspondence may remain uncertain under large viewpoint changes, occlusion, weak target appearance, or inaccurate detections.

Second, multimodal fusion and behavior analysis depend on the coverage and quality of upstream trajectories. When all modalities are severely incomplete or entity embeddings are weak, graph matching and reliability-aware fusion may still produce uncertain correspondences. In addition, group-level behavior evaluation is limited by the lack of complete anomaly annotations, so some results rely on ranking analysis and qualitative inspection.

Finally, the current system is still a research prototype. It verifies the feasibility of the overall pipeline, but real-time processing, large-scale data management, automatic parameter adaptation, and long-term deployment stability still require further work.

\section{Future Work}
Future work can proceed in three directions. First, optical trajectory completion can be strengthened by combining scene-level VPR retrieval with more explicit target-level matching. Object appearance, local geometry, temporal consistency, and detection confidence could be jointly used after Top-\(k\) retrieval.

Second, more complete multimodal datasets are needed. Such datasets should include cross-modal target identities, incomplete trajectory segments, coordinate-missing optical images, and behavior anomaly annotations. They would support more direct evaluation of optical completion, entity registration, fusion quality, and behavior analysis within the same benchmark.

Third, the framework can be improved for practical deployment. Future work should propagate calibrated uncertainty across modules, expand behavior analysis with richer semantic and group context, and optimize the system for efficient retrieval, incremental graph matching, streaming anomaly analysis, and stable long-sequence visualization.
